Passenger detection and counting have evolved from handcrafted visual models to end-to-end deep learning frameworks. This section reviews related work in three domains: traditional methods, deep learning–based detection and tracking, and railway-specific crowd analysis.

\subsection{Traditional Methods}
Before deep learning dominated the field, crowd and passenger counting relied on handcrafted features and shallow models. Techniques such as Histogram of Oriented Gradients (HOG) and Hough transform were applied to detect passengers in bus or platform images \cite{30_khan_passenger_2020,73_Huofu}. While computationally efficient, these classical approaches were highly sensitive to illumination, perspective, and occlusion, resulting in poor generalization to crowded or dynamic scenes. Regression-based models that estimated global crowd density from low-level features \cite{zhang2016single,li2018csrnet} improved robustness but still failed to capture fine-grained individual localization.

\subsection{Deep Learning for Crowd Detection and Counting}
The advent of convolutional neural networks (CNNs) enabled the transition from feature engineering to feature learning. Early deep regression networks \cite{zhang2016single,li2018csrnet} predicted pixel-wise density maps, providing accurate crowd estimates in surveillance footage. However, density-based methods lose instance-level information required for multi-object tracking and identity consistency.  
Object detection frameworks such as Faster R-CNN and YOLO series \cite{redmon2016you,bochkovskiy2020yolov4,74_yolo11_Jocher_Ultralytics_YOLO_2023} addressed this issue by enabling explicit localization. They have since been adapted to head detection tasks \cite{shao2018crowdhuman}, which are more robust in dense or occluded environments. Recent transformer-based detectors further improved long-range context modeling but typically require large datasets and high computational cost, limiting deployment in real-time transportation monitoring.

\subsection{Multi-Object Tracking and Re-Identification}
Multi-object tracking (MOT) extends detection into temporal association. Early online trackers combine Kalman filtering with Hungarian matching, achieving real-time performance\cite{220}. Later works introduced appearance models and motion prediction enhancements, reducing identity switches in crowded scenarios\cite{211}. Appearance-based re-identification (Re-ID) has become essential for maintaining consistent identities across frames\cite{221_OccludedPersonReIdentification_ning2023}. Networks based on EfficientNet and ResNet backbones extract discriminative embeddings that support robust association even under partial occlusions\cite{222,223_Luo_2020}. Despite these advances, most trackers assume a fixed camera and constant-velocity motion model, which breaks down in scenes with strong perspective variation or moving cameras—typical conditions for railway surveillance.

\subsection{Crowd Counting in Transportation Scenarios}
Compared with generic crowd scenes, railway platforms present unique challenges: frequent partial occlusions, fast camera motion, and pronounced scale variation. While vehicle- or station-level passenger analytics have been studied \cite{di_gennaro_hum-card_2024,OpenSensorDataforRail2023}, most approaches rely on full-body detection, making them susceptible to mutual occlusion.  
Recent studies have highlighted the advantages of head-based counting \cite{shao2018crowdhuman}, yet publicly available datasets remain limited and lack domain diversity for railway applications. The absence of unified benchmarks and physically grounded motion models further constrains reproducibility and cross-domain generalization.

\subsection{Summary and Research Gap}
In summary, while deep detection and tracking frameworks have achieved impressive accuracy in general crowd scenes, their performance degrades under railway-specific constraints: nonstationary cameras, narrow viewing angles, and dense human occlusion. Traditional constant-velocity Kalman filters fail to model perspective geometry, causing instability in trajectory estimation and erroneous counts. Moreover, few works explicitly address robust, unique counting—distinguishing repeated entries and exits—from a physically constrained motion perspective.  
This motivates our work: an integrated detect–track–count system that combines head-based detection, appearance-aware re-identification, and a novel physics-constrained Kalman model designed to handle perspective distortion and physically plausible motion in real-time railway surveillance.


